% ============================================
% DOC-04 (EN)
% English translation — SPEC-935-R355
% ============================================
\section*{\textbf{DOC-04} --- Session Report I}
\markboth{DOC-04 --- Session Report I}{}

\noindent \textit{Dr. Heitor Oliveira --- March 14, 1979}

\noindent \textbf{Dr. Heitor Oliveira --- March 14, 1979}

\vspace{0.5em}\noindent\rule{\textwidth}{0.5pt}\vspace{0.5em}

\noindent \textbf{Patient:} Joaquim Antônio Correia --- Patient 1,260
\textbf{Type:} Initial psychiatric assessment session
\textbf{Duration:} 45 minutes

\vspace{0.5em}\noindent\rule{\textwidth}{0.5pt}\vspace{0.5em}

\noindent \textbf{Patient Observations:}

\noindent Joaquim is a thin man, short in stature, with wrinkled skin and a fixed stare. The right eye differs significantly from the left --- it presents an unusual egg-yolk yellow coloration. When questioned, he stated that \textquotedblleft{}a creature\textquotedblright{} changed the color. He refused to elaborate.

\noindent The patient is articulate and intelligent --- far more than the medical record suggests. He demonstrates capacity for metacognition and psychological analysis (e.g., he classified doctors by smell, distinguished between new and experienced doctors). He does not present the classic signs of cognitive deterioration expected of someone with 62 years of institutionalization.

\noindent His speech is lucid, but monofocused on the theme of \textquotedblleft{}hunger\textquotedblright{} and the \textquotedblleft{}creature\textquotedblright{}. When questioned about the homicides, he answered naturally: \textquotedblleft{}hunger killed for me\textquotedblright{\textit{. He shows no remorse, but also no pride --- he shows }acceptance}, as if the events were inevitable.

\noindent \textit{Personal note (manuscript, in the margin):} \textquotedblleft{}The key to this case is not in the delusion. It is in what came before --- the Sertão, the drought, the dead father, the trench. The delusion is only the lock. The story is the key. --- H.O.\textquotedblright{}

\noindent \textbf{Diagnostic Impression (Preliminary):}

\noindent The patient presents a clinical picture compatible with Paranoid Schizophrenia (ICD-8 295.3), with a systematized delusion involving an entity that inhabits him. The structure of the delusion is complex and internally coherent --- unusual for patients institutionalized for decades, who generally present fragmented speech.

\noindent However, there are elements suggesting preserved insight capacity. When I asked whether he thought the creature was real or a product of his mind, he answered: \textit{\textquotedblleft{}Do you think the difference matters?\textquotedblright{}}

\noindent \textbf{Therapeutic Plan:}

\noindent - Weekly assessment sessions
- Request access to the patient's diary (information: the patient has kept a diary for decades; has never allowed anyone to read it)
- Maintain current medication (Chlorpromazine 100mg/day)

\noindent \textbf{Personal Note (not included in the official record):}

\noindent There is something about this patient that unsettles me. I cannot identify it precisely. But at the end of the session, when I stood up to leave, he said: \textit{\textquotedblleft{}You will dream about me tonight.\textquotedblright{}}

\noindent I did dream.

\vspace{0.5em}\noindent\rule{\textwidth}{0.5pt}\vspace{0.5em}

\noindent \textit{Dr. Heitor Oliveira}
\textit{Psychiatrist --- Maximum Security Ward}
\textit{Barbacena Hospital-Colony}

\vspace{0.5em}
\noindent\textit{\textbf{↪ Routes:} \rota{MEM-13} \rotasep \rota{DOC-05} \rotasep \rota{MEM-14} \rotasep \rota{LUC-01}}
